% Gabarit pandoc des corrigés publiés (brut et guidé). Contenu sans marque ni nom de site.
% Compilation : lualatex (latexmk -lualatex).
% Le Markdown reste la source unique : ne pas éditer ce fichier .tex à la main.
\documentclass[11pt,a4paper]{article}

% --- Moteur, polices (Latin Modern) et langue ---------------------------------
\usepackage{iftex}
\ifLuaTeX\else\errmessage{Ce document se compile avec lualatex}\fi
\usepackage{amsmath,amssymb,mathtools}
\usepackage{fontspec}
% mathrm=sym : \mathrm prend ses glyphes dans la police mathématique (ℓ, μ, Ω), comme KaTeX.
\usepackage[math-style=TeX,bold-style=TeX,mathrm=sym]{unicode-math}
% Repli pour les glyphes absents de Latin Modern (indices et exposants Unicode,
% flèches, filets des schémas en texte, symboles) : DejaVu, puis Latin Modern Math
% (ℓ, lettres mathématiques Unicode des légendes).
\directlua{
  luaotfload.add_fallback("repliTexte", {"DejaVuSerif:mode=harf;", "DejaVuSans:mode=harf;", "file:latinmodern-math.otf:mode=harf;"})
  luaotfload.add_fallback("repliMono", {"DejaVuSansMono:mode=harf;", "DejaVuSans:mode=harf;", "file:latinmodern-math.otf:mode=harf;"})
}
\setmainfont{Latin Modern Roman}[RawFeature={fallback=repliTexte}]
\setsansfont{Latin Modern Sans}[RawFeature={fallback=repliTexte}]
\setmonofont{Latin Modern Mono}[RawFeature={fallback=repliMono}]
\setmathfont{Latin Modern Math}
\usepackage[french]{babel}
\frenchsetup{SmallCapsFigTabCaptions=false}

% --- Mise en page --------------------------------------------------------------
\usepackage[a4paper,margin=2.2cm,headheight=14pt,footskip=1.2cm]{geometry}
\usepackage{microtype}
\setlength{\emergencystretch}{4em}
\tolerance=1500
\setlength{\parindent}{0pt}
\setlength{\parskip}{0.55\baselineskip plus 2pt minus 1pt}
\usepackage{xcolor}
% Gris nommés : babel-french rend « ! » actif dans le document.
\definecolor{gris2}{gray}{0.98}
\definecolor{gris3}{gray}{0.97}
\definecolor{gris5}{gray}{0.95}
\definecolor{gris35}{gray}{0.65}
\definecolor{gris40}{gray}{0.60}
\definecolor{gris45}{gray}{0.55}
\definecolor{gris50}{gray}{0.50}
\definecolor{gris55}{gray}{0.45}
\definecolor{gris60}{gray}{0.40}
\definecolor{gris65}{gray}{0.35}
\definecolor{gris70}{gray}{0.30}
\definecolor{gris80}{gray}{0.20}
\definecolor{gris85}{gray}{0.15}
\usepackage{graphicx}
\usepackage{float}
\usepackage{caption}
\captionsetup{font=small,labelformat=empty,justification=centering}
\usepackage{booktabs,longtable,array,multirow,calc}
\usepackage{siunitx}
\sisetup{locale=FR}
\usepackage[version=4]{mhchem}
\usepackage{listings}
\usepackage{fancyhdr}
\usepackage{lastpage}
\usepackage[most]{tcolorbox}
\usepackage{enumitem}
\setlist{nosep,leftmargin=*,topsep=2pt}
\usepackage[hidelinks,unicode]{hyperref}
\usepackage{bookmark}
\hypersetup{pdftitle={Centrale-Supélec TSI — Anglais 2026},pdfauthor={Rédigé par IA},pdfcreator={LaTeX},pdfsubject={Corrigé guidé}}
\usepackage{xurl}

% --- Pandoc : listes, tableaux, figures ----------------------------------------
\providecommand{\passthrough}[1]{#1}
\providecommand{\tightlist}{\setlength{\itemsep}{0pt}\setlength{\parskip}{0pt}}

\makeatletter
\def\maxwidth{\ifdim\Gin@nat@width>\linewidth\linewidth\else\Gin@nat@width\fi}
\def\maxheight{\ifdim\Gin@nat@height>0.55\textheight 0.55\textheight\else\Gin@nat@height\fi}
\makeatother
\setkeys{Gin}{width=\maxwidth,height=\maxheight,keepaspectratio}
\setcounter{secnumdepth}{-\maxdimen}
\providecommand{\pandocbounded}[1]{#1}
\usepackage{longtable,booktabs,array}

% --- Code (listings) -----------------------------------------------------------
\lstset{
  basicstyle=\ttfamily\small,
  columns=fullflexible,
  keepspaces=true,
  breaklines=true,
  breakatwhitespace=false,
  postbreak=\mbox{\textcolor{gris50}{$\hookrightarrow$}\space},
  frame=single,
  rulecolor=\color{gris35},
  framesep=4pt,
  xleftmargin=5pt,xrightmargin=5pt,
  backgroundcolor=\color{gris3},
  keywordstyle=\bfseries,
  commentstyle=\color{gris65},
  stringstyle=\color{gris80},
  showstringspaces=false,
  tabsize=4,
  extendedchars=true,
  aboveskip=6pt,belowskip=6pt,
}

% --- Macros des corrigés (KaTeX ↔ LaTeX) ---------------------------------------
% Les corrigés sont écrits pour KaTeX : on définit ici ce que KaTeX connaît et
% que LaTeX n'a pas, sans jamais retoucher le Markdown.
\providecommand{\R}{\mathbb{R}}
\providecommand{\N}{\mathbb{N}}
\providecommand{\Z}{\mathbb{Z}}
\providecommand{\Q}{\mathbb{Q}}
\providecommand{\C}{\mathbb{C}}
\providecommand{\dd}{\mathrm{d}}
\providecommand{\llbracket}{\lBrack}
\providecommand{\rrbracket}{\rBrack}
\providecommand{\lBrace}{\lbrace}
\providecommand{\rBrace}{\rbrace}
\providecommand{\Box}{\square}
\providecommand{\Diamond}{\lozenge}
\providecommand{\varnothing}{\emptyset}
\providecommand{\argmax}{\operatorname*{arg\,max}}
\providecommand{\argmin}{\operatorname*{arg\,min}}
\renewcommand{\leq}{\leqslant}
\renewcommand{\geq}{\geqslant}

% Équation hors texte : réduite à la largeur de la ligne si elle déborde.
\newsavebox{\boiteequation}
\newcommand{\equationajustee}[1]{%
  \sbox{\boiteequation}{$\displaystyle #1$}%
  \ifdim\wd\boiteequation>\linewidth
    \[\resizebox{\linewidth}{!}{\usebox{\boiteequation}}\]%
  \else
    \[#1\]%
  \fi}

% --- Encadrés pédagogiques (sobres, lisibles en noir et blanc) -----------------
\tcbset{
  encadre/.style={
    enhanced,breakable,sharp corners,
    colback=white,colframe=gris55,
    boxrule=0.4pt,leftrule=2.2pt,
    left=6pt,right=6pt,top=3pt,bottom=3pt,
    before skip=6pt,after skip=6pt,
    fonttitle=\bfseries\small,coltitle=black,
    attach title to upper,after title={\quad},
    parbox=false,
  }
}
\newtcolorbox{indication}[1]{encadre,colframe=gris45,colback=gris2,title={Indication~#1}}
\newtcolorbox{methode}{encadre,colframe=gris70,title={Méthode}}
\newtcolorbox{piege}{encadre,colframe=gris85,boxrule=0.8pt,leftrule=2.2pt,title={Piège}}
\newtcolorbox{jury}{encadre,colframe=gris60,colback=gris5,title={Jury}}
\newtcolorbox{encadrenote}{encadre,colframe=gris40}

% --- Titres de section ---------------------------------------------------------
\usepackage{titlesec}
\titleformat{\section}{\Large\bfseries}{}{0pt}{}
\titleformat{\subsection}{\large\bfseries}{}{0pt}{}
\titleformat{\subsubsection}{\normalsize\bfseries}{}{0pt}{}
\titlespacing*{\section}{0pt}{14pt plus 4pt minus 2pt}{6pt}
\titlespacing*{\subsection}{0pt}{10pt plus 3pt minus 2pt}{4pt}
\titlespacing*{\subsubsection}{0pt}{8pt plus 2pt minus 2pt}{2pt}

% --- En-tête et pied de page ---------------------------------------------------
\pagestyle{fancy}
\fancyhf{}
\fancyhead[L]{\small\textsc{Corrigé guidé} · Centrale-Supélec TSI}
\fancyhead[R]{\small Anglais 2026}
\fancyfoot[L]{\footnotesize Corrigé rédigé par IA, non officiel}
\fancyfoot[R]{\footnotesize \thepage\,/\,\pageref*{LastPage}}
\renewcommand{\headrulewidth}{0.4pt}
\renewcommand{\footrulewidth}{0pt}
\fancypagestyle{plain}{\fancyhf{}\fancyfoot[R]{\footnotesize \thepage\,/\,\pageref*{LastPage}}\fancyfoot[L]{\footnotesize Corrigé rédigé par IA, non officiel}\renewcommand{\headrulewidth}{0pt}}

\begin{document}
\thispagestyle{plain}
\begin{center}
  {\small\scshape Corrigé guidé}\par\smallskip
  {\LARGE\bfseries Centrale-Supélec TSI — Anglais 2026\par}\medskip
  {\large Centrale-Supélec · TSI · Anglais · 2026\par}
\end{center}
\vspace{2pt}\hrule\vspace{8pt}

\emph{Corrigé rédigé par IA, non officiel. Signalez toute erreur.}

\hypertarget{pruxe9sentation-du-sujet}{%
\section{Présentation du sujet}\label{pruxe9sentation-du-sujet}}

\textbf{Format.} L'épreuve dure 4 heures ; calculatrice et tout système électronique sont interdits. Elle comporte un seul exercice :

\begin{longtable}[]{@{}
  >{\raggedright\arraybackslash}p{(\columnwidth - 6\tabcolsep) * \real{0.1101}}
  >{\raggedright\arraybackslash}p{(\columnwidth - 6\tabcolsep) * \real{0.2431}}
  >{\raggedright\arraybackslash}p{(\columnwidth - 6\tabcolsep) * \real{0.2752}}
  >{\raggedright\arraybackslash}p{(\columnwidth - 6\tabcolsep) * \real{0.3416}}@{}}
\toprule\noalign{}
\begin{minipage}[b]{\linewidth}\raggedright
Exercice
\end{minipage} & \begin{minipage}[b]{\linewidth}\raggedright
Nature
\end{minipage} & \begin{minipage}[b]{\linewidth}\raggedright
Longueur
\end{minipage} & \begin{minipage}[b]{\linewidth}\raggedright
Contraintes explicites du sujet
\end{minipage} \\
\midrule\noalign{}
\endhead
\bottomrule\noalign{}
\endlastfoot
Synthèse & synthèse en anglais de quatre documents & 400 mots ± 10 \%, soit 360 à 440 mots, titre inclus & titre obligatoire ; nombre de mots indiqué avec précision à la fin du travail \\
\end{longtable}

Le sujet précise que l'ordre des documents est arbitraire et ne revêt aucune signification particulière. Le barème n'est pas publié dans le sujet. Le fichier du sujet contient aussi, après les quatre pages d'anglais, le sujet « Langue vivante » des autres langues (allemand, arabe, espagnol, italien, portugais, russe, chinois) : ce corrigé ne traite que l'anglais, conformément à l'intitulé de l'épreuve.

\textbf{Pas de rapport de jury.} Centrale-Supélec n'a pas publié de rapport de jury pour la session 2026 au moment de la rédaction. Ce corrigé ne cite donc aucune attente du jury : les conseils qui suivent sont des \textbf{conseils méthodologiques généraux} sur l'exercice de synthèse, présentés comme tels, et non des exigences propres aux correcteurs de cette session.

\textbf{Le dossier.}

\begin{longtable}[]{@{}
  >{\raggedright\arraybackslash}p{(\columnwidth - 6\tabcolsep) * \real{0.1212}}
  >{\raggedright\arraybackslash}p{(\columnwidth - 6\tabcolsep) * \real{0.1366}}
  >{\raggedright\arraybackslash}p{(\columnwidth - 6\tabcolsep) * \real{0.5656}}
  >{\raggedright\arraybackslash}p{(\columnwidth - 6\tabcolsep) * \real{0.1466}}@{}}
\toprule\noalign{}
\begin{minipage}[b]{\linewidth}\raggedright
Document
\end{minipage} & \begin{minipage}[b]{\linewidth}\raggedright
Nature, date
\end{minipage} & \begin{minipage}[b]{\linewidth}\raggedright
Contenu
\end{minipage} & \begin{minipage}[b]{\linewidth}\raggedright
Apport
\end{minipage} \\
\midrule\noalign{}
\endhead
\bottomrule\noalign{}
\endlastfoot
R. Horowitch, \emph{The Atlantic} & enquête, 01/10/2024 & les étudiants des universités d'élite (Columbia, Princeton, Georgetown\ldots) ne parviennent plus à lire des livres entiers ; ils savent déchiffrer mais ne soutiennent plus l'attention ; causes : smartphones (« Being bored has become unnatural »), et un déplacement des valeurs vers l'emploi ; en 1976, 40 \% des lycéens de terminale avaient lu au moins six livres pour le plaisir dans l'année, 11,5 \% aucun, proportions inversées en 2022 ; Columbia allège sa liste de lectures ; la lecture profonde (Maryanne Wolf) développe esprit critique, réflexion sur soi et empathie & le constat et les enjeux \\
J. Rothman, \emph{The New Yorker} & essai, 17/06/2025 & la nature même de la lecture change : « hyper-literacy » pour certains, lecture fragmentée, livre audio, fil d'actualités pour d'autres ; adultes lisant au moins un livre par an : de 55 \% à 48 \% en dix ans ; adolescents de 13 ans lisant pour le plaisir presque tous les jours : de 27 \% à 14 \% ; fin de la « Gutenberg Parenthesis », « secondary orality » (Walter Ong), podcasts, mèmes & une mutation plus qu'une disparition \\
B. Elton, \emph{Blind Faith} & roman dystopique, 2007 & dans une société où lire des romans et penser par soi-même est interdit, Trafford révèle à Sandra Dee un exemplaire des \emph{Hauts de Hurlevent} caché sous la couverture d'un pamphlet ; le livre est « a door » vers « a thousand universes », une évasion par « reason and imagination » & la lecture comme liberté \\
DPM, publicité & image publicitaire, novembre 2021 & « Metaverse --- It's just the beginning » : un homme coiffé d'un casque de réalité virtuelle, allongé dans un fauteuil, dans une pièce posée au milieu d'un paysage préhistorique où rôdent un tyrannosaure et des ptérosaures & l'immersion promise par la technologie \\
\end{longtable}

\textbf{Le thème.} Les deux articles décrivent le recul de la lecture longue chez les jeunes et en cherchent les causes ; le roman et la publicité opposent deux manières de « quitter le monde » : par le livre, qui exige l'imagination du lecteur, ou par la réalité virtuelle, qui la fournit toute faite. La synthèse doit tenir ensemble le constat chiffré, les causes, la nuance de Rothman (la lecture change de forme) et l'enjeu (attention, empathie, imagination).

\textbf{Ce qu'il ne faut pas laisser passer} (conseils généraux) :

\begin{itemize}
\tightlist
\item
  les deux contraintes explicites : un \textbf{titre} et le \textbf{nombre de mots} à la fin, dans la fourchette 360-440 ;
\item
  les \textbf{quatre documents}, y compris la publicité et l'extrait de roman, qui ne sont pas de simples illustrations ;
\item
  la confrontation des documents dans chaque partie, plutôt qu'une suite de résumés ;
\item
  pas d'opinion personnelle ni d'information extérieure au dossier ;
\item
  le statut des documents : l'extrait d'Elton est une fiction, la publicité un discours commercial ; on ne les présente pas comme des faits.
\end{itemize}

\textbf{Gestion du temps} (proposition) : 1 h 15 de lecture et d'analyse, 30 min pour la confrontation et le plan, 1 h 45 de rédaction au brouillon puis au propre, 30 min de décompte et de relecture.

\hypertarget{partie-i-synthuxe8se-de-documents}{%
\section{Partie I --- Synthèse de documents}\label{partie-i-synthuxe8se-de-documents}}

L'unique exercice est découpé en six étapes, de la lecture au décompte final. La solution complète est donnée à l'étape 5.

\hypertarget{uxe9tape-1-lire-le-dossier-et-analyser-chaque-document}{%
\subsection{Étape 1 --- Lire le dossier et analyser chaque document}\label{uxe9tape-1-lire-le-dossier-et-analyser-chaque-document}}

\begin{indication}{1}

Pour chaque document, notez la nature, la date, la source et l'idée principale en une phrase. Relevez les chiffres, mais seulement les plus parlants.

\end{indication}

\begin{indication}{2}

Décrivez la publicité avant de l'interpréter : où est l'homme, que porte-t-il, qu'y a-t-il autour de la pièce, que dit le slogan ?

\end{indication}

\begin{indication}{3}

Dans l'extrait d'Elton, demandez-vous ce que représente le livre caché pour Trafford et ce qu'il oppose à la société qu'il décrit.

\end{indication}

\hypertarget{solution}{%
\subsubsection{Solution}\label{solution}}

Le tableau de la présentation donne l'analyse de chaque document. Les liens à exploiter sont les suivants :

\begin{longtable}[]{@{}
  >{\raggedright\arraybackslash}p{(\columnwidth - 4\tabcolsep) * \real{0.2442}}
  >{\raggedright\arraybackslash}p{(\columnwidth - 4\tabcolsep) * \real{0.5131}}
  >{\raggedright\arraybackslash}p{(\columnwidth - 4\tabcolsep) * \real{0.2127}}@{}}
\toprule\noalign{}
\begin{minipage}[b]{\linewidth}\raggedright
Lien
\end{minipage} & \begin{minipage}[b]{\linewidth}\raggedright
Documents
\end{minipage} & \begin{minipage}[b]{\linewidth}\raggedright
Usage dans la synthèse
\end{minipage} \\
\midrule\noalign{}
\endhead
\bottomrule\noalign{}
\endlastfoot
recul chiffré de la lecture & Horowitch, Rothman & partie I \\
écrans et attention & Horowitch (smartphones, TikTok), Rothman (YouTube, Netflix), Elton (« infotainment loops », écrans sur tous les murs), publicité & partie II \\
la lecture change de forme & Rothman (« hyper-literacy », audio, oralité seconde) & nuance de la partie II \\
s'évader dans un autre monde & Elton (le livre comme porte), publicité (le métavers comme paysage) & partie III \\
enjeux de la lecture profonde & Horowitch (empathie, esprit critique), Elton (raison, imagination) & partie III \\
\end{longtable}

\begin{piege}

Horowitch distingue soigneusement les étudiants d'élite, qui savent déchiffrer mais n'arrivent plus à se concentrer, de ceux qui ont de vraies difficultés de lecture dans d'autres établissements. Écrire que les étudiants de Columbia « ne savent plus lire » est un contresens.

\end{piege}

\hypertarget{uxe9tape-2-confronter-les-documents-et-formuler-la-probluxe9matique}{%
\subsection{Étape 2 --- Confronter les documents et formuler la problématique}\label{uxe9tape-2-confronter-les-documents-et-formuler-la-probluxe9matique}}

\begin{indication}{1}

Les articles constatent, le roman et la publicité proposent deux formes d'évasion. Quelle question permet de faire dialoguer ces deux groupes ?

\end{indication}

\begin{indication}{2}

Pensez à ce qui disparaît (ou non) quand la lecture longue recule.

\end{indication}

\hypertarget{solution-1}{%
\subsubsection{Solution}\label{solution-1}}

\textbf{Problématique} : que perd-on lorsque la lecture longue et attentive des livres cède la place aux expériences numériques ?

Elle convoque les quatre documents : les articles mesurent le recul, Rothman nuance (la lecture se transforme), Horowitch et Elton disent ce qui est en jeu, la publicité incarne l'alternative numérique.

\begin{methode}

Une problématique de synthèse est une question à laquelle les documents apportent des réponses différentes. Une question à laquelle tous répondent de la même façon ne produit qu'un résumé.

\end{methode}

\hypertarget{uxe9tape-3-construire-le-plan}{%
\subsection{Étape 3 --- Construire le plan}\label{uxe9tape-3-construire-le-plan}}

\begin{indication}{1}

Trois parties pour 400 mots : un constat, ses causes et ses nuances, puis l'enjeu.

\end{indication}

\begin{indication}{2}

Gardez le roman et la publicité pour les parties où ils éclairent le mieux le propos : ils ne prouvent rien chiffre à l'appui, mais ils montrent deux imaginaires de l'évasion.

\end{indication}

\hypertarget{solution-2}{%
\subsubsection{Solution}\label{solution-2}}

\begin{longtable}[]{@{}
  >{\raggedright\arraybackslash}p{(\columnwidth - 4\tabcolsep) * \real{0.1059}}
  >{\raggedright\arraybackslash}p{(\columnwidth - 4\tabcolsep) * \real{0.3617}}
  >{\raggedright\arraybackslash}p{(\columnwidth - 4\tabcolsep) * \real{0.5024}}@{}}
\toprule\noalign{}
\begin{minipage}[b]{\linewidth}\raggedright
Partie
\end{minipage} & \begin{minipage}[b]{\linewidth}\raggedright
Idée directrice
\end{minipage} & \begin{minipage}[b]{\linewidth}\raggedright
Documents mobilisés
\end{minipage} \\
\midrule\noalign{}
\endhead
\bottomrule\noalign{}
\endlastfoot
I & Un recul net de la lecture longue, jusque dans les universités d'élite & Rothman (55 → 48 \%, 27 → 14 \%), Horowitch (Columbia, listes de lectures allégées) \\
II & Des causes multiples, et une lecture qui change de forme & Horowitch (smartphones, ennui, choix de carrière), Rothman (« hyper-literacy », audio, oralité), publicité (immersion technologique) \\
III & Ce que les livres offrent et que les écrans ne remplacent pas & Horowitch (lecture profonde, empathie), Elton (le livre comme porte vers mille univers), publicité (évasion fournie par la technologie) \\
\end{longtable}

Répartition visée : introduction ≈ 75 mots, parties ≈ 85, 125 et 95 mots, conclusion ≈ 45 mots.

\begin{piege}

Ne réservez pas un paragraphe entier à la publicité ou au roman : ils doivent dialoguer avec les articles.

\end{piege}

\hypertarget{uxe9tape-4-choisir-le-titre-et-ruxe9diger-lintroduction}{%
\subsection{Étape 4 --- Choisir le titre et rédiger l'introduction}\label{uxe9tape-4-choisir-le-titre-et-ruxe9diger-lintroduction}}

\begin{indication}{1}

Le titre est obligatoire et compté. Un jeu de mots sur le livre ou la page fonctionne bien s'il reste clair.

\end{indication}

\begin{indication}{2}

L'introduction présente les quatre documents (auteur ou source, date, nature) et se termine par la problématique.

\end{indication}

\hypertarget{solution-3}{%
\subsubsection{Solution}\label{solution-3}}

Titre retenu : \textbf{Turning the Page on Books?} (5 mots). Autres possibilités : \emph{Reading in the Age of Screens} ; \emph{Closing the Book?}

In a world of screens, is reading books becoming a thing of the past? Two American articles, by Rose Horowitch in The Atlantic and Joshua Rothman in The New Yorker, examine the decline of reading, while an extract from Ben Elton's dystopian novel Blind Faith and a 2021 advertisement for the Metaverse offer two contrasting visions of escape. They raise the question of what is lost when sustained reading gives way to digital experiences.

(74 mots.)

\begin{methode}

Présentez les documents par groupes (les deux articles, puis le roman et la publicité) : c'est plus court qu'une énumération et cela annonce déjà la confrontation.

\end{methode}

\hypertarget{uxe9tape-5-ruxe9diger-la-synthuxe8se-compluxe8te}{%
\subsection{Étape 5 --- Rédiger la synthèse complète}\label{uxe9tape-5-ruxe9diger-la-synthuxe8se-compluxe8te}}

\begin{indication}{1}

Chaque partie s'ouvre par une phrase-thèse (\emph{First, the articles describe\ldots{}}, \emph{However, the documents suggest\ldots{}}), puis fait parler deux ou trois documents.

\end{indication}

\begin{indication}{2}

Choisissez deux ou trois chiffres seulement et donnez-en le sens (\emph{roughly halved}) plutôt que de les aligner.

\end{indication}

\hypertarget{solution-4}{%
\subsubsection{Solution}\label{solution-4}}

\textbf{Turning the Page on Books?}

In a world of screens, is reading books becoming a thing of the past? Two American articles, by Rose Horowitch in The Atlantic and Joshua Rothman in The New Yorker, examine the decline of reading, while an extract from Ben Elton's dystopian novel Blind Faith and a 2021 advertisement for the Metaverse offer two contrasting visions of escape. They raise the question of what is lost when sustained reading gives way to digital experiences.

First, the articles describe a striking decline. Rothman notes that the share of American adults who read at least one book a year fell from 55 to 48 percent in a decade, while the proportion of thirteen-year-olds reading for fun almost every day was roughly halved. Horowitch adds that even students at elite universities such as Columbia, who can decode texts perfectly well, now feel overwhelmed by the idea of reading several books a semester, so that some professors have shortened their reading lists.

Both journalists blame screens for this change. Smartphones, TikTok or Netflix constantly compete for attention, and, as a psychologist quoted by Horowitch puts it, ``being bored has become unnatural''. Yet the problem is not only technical: many students choose not to read because they care more about their future careers. Rothman also insists that reading is changing rather than disappearing: some people enjoy a kind of ``hyper-literacy'', others switch between e-books and audio narration, and theorists even speak of a return to an oral culture of comments, podcasts and memes. The Metaverse advertisement goes further, promising total immersion: a man wearing a headset lies back in his living room while dinosaurs roam the wild landscape around him, under the slogan ``It's just the beginning''.

However, the documents suggest that books offer something screens cannot replace. For Horowitch, ``deep reading'' develops critical thinking, self-reflection and empathy, by transporting readers into other minds, but only if they read whole works rather than excerpts. Elton's novel turns this idea into fiction: in a society where private thought is forbidden, Trafford sees a hidden copy of Wuthering Heights as ``a door'' to ``a thousand universes'', a way to escape through ``reason and imagination''. The advertisement promises a similar escape, yet one delivered by technology rather than built by the reader's own imagination.

In conclusion, the documents agree that sustained reading is losing ground to digital media, even if reading itself is taking new forms. What is at stake is not just an old habit, but our capacity to concentrate, to imagine and to understand others.

\textbf{Number of words (title included): 424}

\begin{methode}

Les verbes introducteurs varient et situent la position de chaque auteur : \emph{notes}, \emph{adds}, \emph{insists}, \emph{suggest}. Ils évitent la répétition de \emph{says}.

\end{methode}

\begin{piege}

Les chiffres de Rothman et ceux d'Horowitch ne portent pas sur les mêmes populations (adultes et adolescents de 13 ans d'un côté, lycéens de terminale de l'autre). Ne les mélangez pas.

\end{piege}

\hypertarget{uxe9tape-6-compter-les-mots-et-relire}{%
\subsection{Étape 6 --- Compter les mots et relire}\label{uxe9tape-6-compter-les-mots-et-relire}}

\begin{indication}{1}

Comptez paragraphe par paragraphe, titre compris, et reportez le total exact.

\end{indication}

\begin{indication}{2}

Relisez les temps (\emph{fell \ldots{} in a decade}, \emph{was roughly halved}), les articles et l'accord sujet-verbe après \emph{the share of} et \emph{the proportion of}.

\end{indication}

\hypertarget{solution-5}{%
\subsubsection{Solution}\label{solution-5}}

\begin{longtable}[]{@{}ll@{}}
\toprule\noalign{}
Bloc & Mots \\
\midrule\noalign{}
\endhead
\bottomrule\noalign{}
\endlastfoot
Titre & 5 \\
Introduction & 74 \\
Partie I & 84 \\
Partie II & 124 \\
Partie III & 94 \\
Conclusion & 43 \\
\textbf{Total} & \textbf{424} \\
\end{longtable}

Le total (424 mots) est dans la fourchette 360-440, quelle que soit la convention adoptée pour les traits d'union et les apostrophes (voir la revue).

Points de langue vérifiés :

\begin{itemize}
\tightlist
\item
  \emph{the share of American adults \ldots{} fell}, \emph{the proportion of thirteen-year-olds \ldots{} was roughly halved} : verbe au singulier, accordé à \emph{share} et \emph{proportion} ;
\item
  \emph{thirteen-year-olds} : pas de \emph{s} à \emph{year}, trait d'union entre les trois éléments ;
\item
  \emph{care more about their future careers} : \emph{care about}, pas \emph{care of} ;
\item
  \emph{lies back in his living room} : \emph{lie, lay, lain} (s'allonger), à ne pas confondre avec \emph{lay, laid, laid} (poser).
\end{itemize}

\begin{piege}

Un nombre de mots inexact ou absent, ou un total hors fourchette, contrevient à une consigne explicite du sujet. Indiquer « about 400 words » ne répond pas à la demande d'indiquer le nombre « avec précision ».

\end{piege}

\hypertarget{bilan}{%
\section{Bilan}\label{bilan}}

\begin{itemize}
\tightlist
\item
  \textbf{Méthode de la synthèse} : analyse document par document, grille de confrontation, problématique tirée des différences de point de vue, plan en trois parties qui croise les sources, rédaction neutre et attribuée.
\item
  \textbf{Consignes de forme} : titre obligatoire, décompte exact à la fin, fourchette de 360 à 440 mots, titre inclus.
\item
  \textbf{Documents non argumentatifs} : un extrait de roman et une publicité se décrivent puis s'interprètent en lien avec les articles, sans être présentés comme des faits.
\item
  \textbf{Langue utile} : \emph{attention span}, \emph{deep reading}, \emph{to skim}, \emph{to be overwhelmed by}, \emph{reading list}, \emph{literacy}, \emph{audiobook}, \emph{immersive}, \emph{virtual reality headset}, \emph{to escape into}.
\item
  \textbf{Transposable} : les dossiers sur les écrans, l'attention, l'école et la culture numérique se traitent avec la même grille (constat chiffré, causes, nuances, enjeux).
\end{itemize}

\end{document}
